% !TEX root = task-2-8-1.tex
% Editable task bank. Append new BankTask blocks at the end of the document.
\documentclass[12pt,a4paper]{article}
\usepackage{physics-textbook}
\usepackage{graphicx,needspace}
\usetikzlibrary{decorations.pathmorphing}
\newif\ifShowSolutions
\ifdefined\TutorEdition
  \ShowSolutionstrue
\else
  \ShowSolutionsfalse % Change to \ShowSolutionstrue to display solutions.
\fi
\newcounter{banktask}
\newsavebox{\BankTaskBox}
\newcommand{\BankTaskContent}[3]{%
  \noindent\textbf{\thebanktask.}\quad #2\par
  \ifShowSolutions
    \Needspace{50mm}% Keep the solution heading with its opening text or diagram.
    #3\par
    \ifdefined\TutorEdition
      \ifcsname TutorSource#1\endcsname
        \nopagebreak\smallskip{\small\csname TutorSource#1\endcsname\par}%
      \fi
    \fi
  \fi
}
\newcommand{\BankTask}[3]{%
  \refstepcounter{banktask}\par\medskip
  \sbox{\BankTaskBox}{\begin{minipage}{\linewidth}%
    \BankTaskContent{#1}{#2}{#3}\end{minipage}}%
  \ifdim\dimexpr\ht\BankTaskBox+\dp\BankTaskBox\relax>\dimexpr\textheight-20mm\relax
    \Needspace{35mm}\BankTaskContent{#1}{#2}{#3}%
  \else
    \Needspace{\dimexpr\ht\BankTaskBox+\dp\BankTaskBox\relax}%
    \noindent\usebox{\BankTaskBox}\par
  \fi
  \medskip
}
\begin{document}
\BookSubsection{2.8.1}{Внутренняя энергия идеального одноатомного газа}


% Source sdamgia-19661: https://phys-ege.sdamgia.ru/problem?id=19661
\BankTask{sdamgia-19661}{
Порция идеального одноатомного газа при температуре $300\,\text{К}$ обладает внутренней энергией $1{,}242\,\text{Дж}$. Сколько атомов содержит эта порция газа? (В $10^{20}$; целое число.)
}{
\smallskip\noindent\textbf{Решение}\par
Внутренняя энергия идеального одноатомного газа: 

$U = \frac{3}{2} \nu RT = \frac{3}{2} \cdot \frac{N}{N_A} RT.$

Откуда и находим число атомов: 

$N = \frac{2UN_A}{3RT} = 2 \cdot 10 ^{20} .$

Ответ: $2$.
}

% Source sdamgia-3890: https://phys-ege.sdamgia.ru/problem?id=3890
\BankTask{sdamgia-3890}{
Идеальный одноатомный газ, находящийся при температуре $T$, нагрели до температуры $2$ $T$, сообщив ему количество теплоты $10\,\text{Дж}$. В результате гaз совершил работу $5\,\text{Дж}$. Какое количество теплоты отдаст газ, если его после этого изохорически охладить до температуры $\frac{3T}{2} .$ (В джоулях, округлите до десятых.)
}{
\smallskip\noindent\textbf{Решение}\par
Согласно первому началу термодинамики, переданное газу тепло связано с изменением внутренней энергией и работой газа против внешних сил соотношением: $Q=A + \Delta U.$ Для идеального одноатомного газа изменение внутренней энергии определяется только изменением температуры: $\Delta U= \frac{3}{2} \nu R \Delta T .$ 

Таким образом, изменение внутренней энергии в первом процессе равно:

$\Delta U_1= \frac{3}{2} \nu R ( 2T - T ) = \frac{3}{2} \nu R T= Q_1 - A_1=10\text{Дж} - 5\text{Дж}=5 \text{Дж}$

Для второго процесса имеем (работа равна нулю, поскольку процесс идет изохорически):

$Q_2=\Delta U_2= \frac{3}{2} \nu R ( 3T/2 - 2T ) = - \frac{1}{2} \cdot \frac{3}{2} \nu R T= - \frac{1}{2} \cdot 5\text{Дж}= - 2{,}5 \text{Дж}$

Следовательно, при данном изохорическом охлаждении газ отдаст $2{,}5\,\text{Дж}$.

Ответ: $2{,}5\,\text{Дж}$.
}

% Source sdamgia-31881: https://phys-ege.sdamgia.ru/problem?id=31881
\BankTask{sdamgia-31881}{
В сосуде находился идеальный одноатомный газ при температуре $250\,\text{К}$. Половину газа выпустили из сосуда, а оставшийся в нем газ нагрели до температуры $600\,\text{К}$. Во сколько раз внутренняя энергия содержимого сосуда в конечном состоянии отличается от внутренней энергии содержимого сосуда в начальном состоянии?
}{
\smallskip\noindent\textbf{Решение}\par
Внутренняя энергия одноатомного идеального газа в первом состоянии равна $U_1= \frac{3}{2} \nu RT_1,$ во втором состоянии — $U_2= \frac{3}{2} \cdot \frac{\nu}{2} RT_2.$ Тогда найдем, во сколько раз изменилась внутренняя энергия:

$\frac{U_2}{U_1} = \frac{T_2}{2T_1} = \frac{600}{2 \cdot 250} =1{,}2.$

Ответ: $1{,}2$.
}

% Source sdamgia-37890: https://phys-ege.sdamgia.ru/problem?id=37890
\BankTask{sdamgia-37890}{
При сжатии идеального одноатомного газа при постоянном давлении внешние силы совершили работу $800\,\text{Дж}$. Чему равно изменение внутренней энергии газа? (В джоулях.)
}{
\smallskip\noindent\textbf{Решение}\par
Работа газа при изобарном процессе равна $A=\nu R\Delta T= - 800\text{Дж},$ изменение внутренней энергии одноатомного идеального газа при изобарном процессе: 

$\Delta U = \frac{3}{2} \nu R\Delta T = \frac{3}{2} A = \frac{3}{2} \cdot ( - 800 ) = - 1200 \text{Дж}.$

Ответ: $-1200$.
}

% Source sdamgia-11934: https://phys-ege.sdamgia.ru/problem?id=11934
\BankTask{sdamgia-11934}{
Сосуд разделен на две равные по объему части пористой неподвижной перегородкой. В начальный момент времени в левой части сосуда содержится $4\,\text{моль}$ гелия, в правой — $40\,\text{г}$ аргона. Перегородка может пропускать молекулы гелия и является непроницаемой для молекул аргона. Температура газов одинаковая и остается постоянной.

$1$. Концентрация гелия в правой части сосуда в $2$ раза меньше, чем аргона.

$2$. Отношение давления газов в правой части сосуда к давлению газа в левой части равно $1{,}5$.

$3$. В правой части сосуда общее число молекул газов меньше, чем в левой части.

$4$. Внутренняя энергия гелия и аргона одинакова.

$5$. В результате установления равновесия давление в правой части сосуда увеличилось в $3$ раза.
}{
\smallskip\noindent\textbf{Решение}\par
Перегородка проницаема только для молекул гелия, поэтому в результате установления равновесия парциальное давление гелия в левой части будет равно парциальному давлению гелия в правой части. Давление газа можно вычислить по формуле: $p = \frac{\nu RT}{V} .$ Парциальные давления гелия в левой и правой части одинаковы, одинаковы температуры и объемы частей, следовательно, одинаковы и количества вещества гелия в левой и правой частях сосуда, то есть в левой и правой части сосуда будет содержаться по $2$ моля гелия.

Найдем связь концентрации и количества вещества: $n = \frac{N}{V} = \frac{\nu N_{\text{А}}}{V} .$ То есть концентрации и количества вещества зависят прямо пропорционально друг от друга, также заметим, что чем больше количество вещества, тем больше и количество молекул.

Найдем количество вещества аргона: $\nu_{\mathrm{Ar}} = \frac{m_{\mathrm{Ar}}}{\mu_{\mathrm{Ar}}} = \frac{40\text{г}}{40\text{г}/\text{моль}} = 1\text{моль}.$

Используя полученное выше, рассмотрим данные в задании утверждения.

$1$. Концентрация гелия в два раза больше концентрации аргона в правой части сосуда. Утверждение $1$ неверно.

$2$. Найдем отношение давлений газов в правой и левой частях сосуда: $\frac{p_{\text{п}}}{p_{\text{л}}} = \frac{\nu_{\text{г.п.}} + \nu_{\mathrm{Ar}}}{\nu_{\text{г.л.}}} = \frac{2\text{моль} + 1\text{моль}}{2\text{моль}} = 1{,}5,$ здесь $\nu_{\text{г.п.}},\nu_{\text{г.л.}},\nu_{\mathrm{Ar}}$ — количество вещества гелия в правой части, количество вещества гелия в левой части, количество вещества аргона, соответственно. Утверждение $2$ верно.

$3$. Количество вещества газов в правой части сосуда больше количества вещества газа в левой части сосуда, следовательно, в правой части сосуда общее число молекул газа больше, чем в левой части сосуда. Утверждение $3$ неверно.

$4$. Внутренняя энергия одноатомного идеального газа может быть вычислена по формуле: $U = \frac{3}{2} \nu RT.$ Температура газов одинакова. Количество вещества гелия больше количества вещества аргона, следовательно, внутренняя энергия гелия больше внутренней энергии аргона. Утверждение $4$ неверно.

$5$. Найдем отношение конечного давления в правой части сосуда к начальному давлению в правой части сосуда: $\frac{p_{\text{к}}}{p_{\text{н}}} = \frac{\nu_{\text{г.п.}} + \nu_{\mathrm{Ar}}}{\nu_{\mathrm{Ar}}} = \frac{2\text{моль} + 1\text{моль}}{1\text{моль}} = 3.$ Утверждение $5$ верно.

Ответ: $25$.
}

% Source sdamgia-51361: https://phys-ege.sdamgia.ru/problem?id=51361
\BankTask{sdamgia-51361}{
В двух различных сосудах находятся по $1\,\text{моль}$ идеальных газов. Объем первого сосуда в два раза больше, чем второго. В первом сосуде находится неон при температуре $27^\circ\text{C}$, во втором — аргон при температуре $450\,\text{К}$.

$1$. Концентрация неона в два раза больше, чем аргона.

$2$. Среднеквадратичные скорости молекул одинаковы.

$3$. Давление аргона в $3$ раза больше, чем неона. 

$4$. Средняя кинетическая энергия теплового движения молекул в первом сосуде в $2$ раза больше, чем во втором.

$5$. Внутренняя энергия неона в $1{,}5$ раза меньше, внутренней энергии аргона.
}{
\smallskip\noindent\textbf{Решение}\par
$1$) Неверно. Концентрация газа равна $n = \frac{N}{V} = \frac{\nu N_A}{V} .$ Из формулы следует, что при равных количествах вещества объем неона в $2$ раза больше объема аргона, значит, концентрация неона в $2$ раза меньше концентрации аргона.

$2$) Неверно. Средняя квадратичная скорость теплового движения молекул определяется формулой $\overline v = \sqrt{\frac{3RT}{M}} .$ Найдем отношение скоростей для газов:

$\frac{\overline v _{\text{н}}}{\overline v _{\text{а}}} = \sqrt{\frac{T_{\text{н}}M_{\text{а}}}{T_{\text{а}}M_{\text{н}}}} = \sqrt{\frac{300 \cdot 0{,}04}{0{,}02 \cdot 450}} \approx 1{,}15 > 1,$

следовательно, средняя квадратичная скорость молекул неона больше средней квадратичной скорости молекул аргона.

$3$) Верно. Из уравнения Клапейрона — Менделеева $pV = \nu RT$ находим отношение давлений газов:

$\frac{p_{\text{н}}}{p_{\text{а}}} = \frac{T_{\text{н}}V_{\text{а}}}{T_{\text{а}}V_{\text{н}}} = \frac{300 \cdot V_{\text{а}}}{450 \cdot 2V_{\text{а}}} = \frac{1}{3} ,$

следовательно, давление аргона в $3$ раза больше давления неона.

$4$) Неверно. Средняя кинетическая энергия молекул газа $\overline{E}_{\text{к}} = \frac{3}{2} kT.$ Тогда получаем:

$\frac{\overline{E}_{\text{н}}}{\overline{E}_{\text{а}}} = \frac{T_{\text{н}}}{T_{\text{а}}} = \frac{300}{450} = \frac{2}{3} ,$

значит, средняя кинетическая энергия молекул аргона в $1{,}5$ раза больше средней кинетической энергии молекул неона.

$5$) Верно. Внутренняя энергия газа $U = \frac{3}{2} \nu RT.$ Тогда из соотношения

$\frac{U_{\text{а}}}{U_{\text{н}}} = \frac{T_{\text{а}}}{T_{\text{н}}} = \frac{450}{300} = 1{,}5$

следует, что внутренняя энергия аргона в $1{,}5$ раза больше внутренней энергии неона.

Ответ: $35$.
}

% Source sdamgia-43058: https://phys-ege.sdamgia.ru/problem?id=43058
\BankTask{sdamgia-43058}{
В результате охлаждения внутренняя энергия одноатомного идеального газа уменьшилась в $1{,}5$ раза. Чему равна начальная температура газа, если его конечная температура $200\,\text{К}$? (В кельвинах.)
}{
\smallskip\noindent\textbf{Решение}\par
Внутренняя энергия одноатомного идеального газа $U= \frac{3}{2} \nu RT.$ При ее уменьшении в $1{,}5$ раза абсолютная температура уменьшилась в $1{,}5$ раза. Следовательно, начальная температура равнялась $T_1=1{,}5T_2=1{,}5 \cdot 200=300\text{К}.$

Ответ: $300$.
}

% Source sdamgia-9100: https://phys-ege.sdamgia.ru/problem?id=9100
\BankTask{sdamgia-9100}{
В сосуде содержится $0{,}1\,\text{моль}$ аргона. Среднеквадратичная скорость его молекул равна $400\,\text{м/с}$. Чему равна внутренняя энергия этой порции аргона?
}{
\smallskip\noindent\textbf{Решение}\par
Все инертные газы являются одноатомными, аргон — инертный газ, следовательно, одноатомный. Внутренняя энергия одноатомного газа: $U= \frac{3}{2} \nu RT= \frac{3}{2} pV.$ 

Запишем основное уравнение МКТ газа: $p= \frac{1}{3} m_0n\overline v ^{2} = \frac{1}{3} \rho \overline v ^{2} = \frac{1}{3} \frac{m}{V} \overline v ^{2} = \frac{1}{3} \frac{\nu M}{V} \overline v ^{2} .$ 

Найдем внутреннюю энергию аргона:

$U= \frac{3}{2} \cdot \frac{1}{3} \cdot \frac{\nu M}{V} \cdot \overline v ^{2} V= \frac{1}{2} \nu M\overline v ^{2} = \frac{1}{2} \cdot 0{,}1 \cdot 0{,}04 \cdot 400 ^{2} =320\text{Дж}.$

Ответ: $320\,\text{Дж}$.
}

% Source sdamgia-6782: https://phys-ege.sdamgia.ru/problem?id=6782
\BankTask{sdamgia-6782}{
Теплоизолированный сосуд объемом $V=4\,\text{м}^{3}$ разделен пористой перегородкой на две равные части. В начальный момент в одной части сосуда находится $1\,\text{моль}$ гелия, а в другой — $1\,\text{моль}$ неона. Атомы гелия могут свободно проникать через перегородку, а атомы неона — нет. Начальная температура гелия равна температуре неона: $T=400\,\text{К}$. Определите внутреннюю энергию газа в той части сосуда, где первоначально находился неон, после установления равновесия в системе.
}{
\smallskip\noindent\textbf{Решение}\par
По условию задачи при установлении термодинамического равновесия газ в сосуде не обменивается теплом с внешним миром и не совершает работы. Поэтому внутренняя энергия газа сохраняется, вследствие чего температура газа в конечном состоянии такая же, как и в начальном. После установления равновесия в системе гелий равномерно распределится по всему сосуду, температура газов не изменится. В результате количество газа в той части сосуда, где первоначально находился неон, окажется $\nu_1=\nu/2 = 0{,}5\text{моль}$ гелия и $\nu_2=\nu=1\text{моль}$ неона. Внутренняя энергия газа пропорциональна температуре и количеству молей вещества. Следовательно, внутренняя энергия смеси газов равна: $U = \frac{3}{2} \nu RT,$ где $\nu=\nu_1 + \nu_2.$ Окончательно получим: $U= \frac{3}{2} ( \frac{\nu}{2} + \nu ) RT= \frac{9}{4} \nu RT= \frac{9}{4} \cdot 1 \cdot 8{,}31 \cdot 400\approx 7480\text{Дж}.$

Ответ: $7480\,\text{Дж}$.
}

% Source sdamgia-23248: https://phys-ege.sdamgia.ru/problem?id=23248
\BankTask{sdamgia-23248}{
В теплоизолированном сосуде с жесткими стенками находятся $0{,}1$ моля идеального одноатомного газа и пружинный маятник. Жесткость пружины маятника $1000\,\text{Н/м}$, амплитуда колебаний его груза $10\,\text{см}$. Считая, что нагревается только газ, найдите, на сколько градусов повысится температура газа после того, как колебания маятника прекратятся из-за действия силы вязкого трения. (До целого числа.)
}{
\smallskip\noindent\textbf{Решение}\par
Колеблющийся маятник обладает энергией $W= \frac{kA ^{2}}{2} .$ Колебания являются затухающими, потому что из-за вязкого трения данная энергия передается в виде теплоты газу. Поскольку газ находится в жестком сосуде, то его объем неизменен, поэтому работу газ не совершает. Тогда вся полученная теплота идет на увеличение внутренней энергии газа $\Delta U= \frac{3}{2} \nu R\Delta T.$ Отсюда находим изменение температуры газа: $\Delta T= \frac{kA ^{2}}{3\nu R} = \frac{1000 \cdot 0{,}01}{3 \cdot 0{,}1 \cdot 8{,}31} \approx 4\text{К}.$

Ответ: $4\,\text{К}$.
}


% Source 90F3B1: https://ege.fipi.ru/bank/questions.php?search=1&proj=BA1F39653304A5B041B656915DC36B38&qid=90F3B1
\BankTask{90F3B1}{
При изобарном расширении постоянной массы гелия при нормальном атмосферном давлении его объём увеличился на $0{,}2\,\text{м}^3$. Определите изменение внутренней энергии гелия. (в $\text{кДж}$).
}{
\smallskip\noindent\textbf{Решение от GPT}\par
Для одноатомного идеального газа $U=\frac32pV$. При постоянном давлении $\Delta U=\frac32p\Delta V=\frac32\cdot10^5\cdot0{,}2=3\cdot10^4\,\text{Дж}=30\,\text{кДж}$. Ответ: $30\,\text{кДж}$.
}

% Source EE7CA0: https://ege.fipi.ru/bank/questions.php?search=1&proj=BA1F39653304A5B041B656915DC36B38&qid=EE7CA0
\BankTask{EE7CA0}{
Идеальный одноатомный газ находится в сосуде с жёсткими стенками объёмом \(0{,}6\ \text{м}^{3}\). При охлаждении его давление снизилось на \(3\cdot10^{3}\ \text{Па}\).

На сколько уменьшилась внутренняя энергия газа (в $\text{кДж}$)?
}{
\smallskip\noindent\textbf{Решение от GPT}\par
Для одноатомного идеального газа \(U=\frac32pV\). При постоянном объёме уменьшение внутренней энергии равно
\[
U_1-U_2=\frac32V(p_1-p_2)=\frac32\cdot0{,}6\cdot3\cdot10^3=2700\ \text{Дж}=2{,}7\ \text{кДж}.
\]
Ответ: \(2{,}7\ \text{кДж}\).
}

% Source 1B1B5D: https://ege.fipi.ru/bank/questions.php?search=1&proj=BA1F39653304A5B041B656915DC36B38&qid=1B1B5D
\BankTask{1B1B5D}{
При изобарном расширении неона его объём увеличился на \(0{,}4\,\text{м}^3\), а внутренняя энергия увеличилась на \(6\,\text{кДж}\). Определите давление неона. (в $\text{кПа}$).
}{
\smallskip\noindent\textbf{Решение от GPT}\par
Для одноатомного идеального газа \(U=\frac32pV\). При постоянном давлении \(\Delta U=\frac32p\Delta V\), откуда
\[
p=\frac{2\Delta U}{3\Delta V}=\frac{2\cdot6\cdot10^3}{3\cdot0{,}4}=10^4\,\text{Па}=10\,\text{кПа}.
\]
Ответ: \(10\,\text{кПа}\).
}

% Source 84BDA5: https://ege.fipi.ru/bank/questions.php?search=1&proj=BA1F39653304A5B041B656915DC36B38&qid=84BDA5
\BankTask{84BDA5}{
Сосуд разделён на две равные по объёму части пористой неподвижной перегородкой. В левой части сосуда содержится $8\ \text{г}$ гелия, в правой --- $1\ \text{моль}$ аргона. Перегородка может пропускать молекулы гелия и является непроницаемой для молекул аргона. Температура газов одинакова и остаётся постоянной.
Рассматривается состояние газов после установления равновесия в системе.
\begin{enumerate}
\item Давление в обеих частях сосуда одинаково.
\item Концентрация гелия и аргона в правой части сосуда одинакова.
\item Внутренняя энергия гелия в сосуде больше, чем внутренняя энергия аргона.
\item Внутренняя энергия гелия в сосуде в конечном состоянии больше, чем в начальном.
\item В правой части сосуда общее число молекул газов в $2$ раза меньше, чем в левой части.
\end{enumerate}
}{
\smallskip\noindent\textbf{Решение от GPT}\par
Молярная масса гелия $M_{\mathrm{He}}=4\ \text{г/моль}$, поэтому $\nu_{\mathrm{He}}=8/4=2\ \text{моль}$. В равновесии концентрация гелия одинакова в обеих частях: в каждой находится по $1\ \text{моль}$ гелия. Весь аргон ($1\ \text{моль}$) остаётся справа. Пусть объём каждой части равен $V$.
\begin{enumerate}
\item Неверно: $p_{\text{лев}}=\dfrac{(1\ \text{моль})RT}{V}$, $p_{\text{прав}}=2p_{\text{лев}}$.
\item Верно: справа количества гелия и аргона равны, поэтому равны и их концентрации.
\item Верно: оба газа одноатомные, $U=\frac32\nu RT$; следовательно, $U_{\mathrm{He}}/U_{\mathrm{Ar}}=2$.
\item Неверно: количество гелия и его температура не изменились, поэтому его внутренняя энергия постоянна.
\item Неверно: справа общее число молекул в $2$ раза больше, чем слева.
\end{enumerate}
Ответ: $23$.
}

% Source 0443DE: https://ege.fipi.ru/bank/questions.php?search=1&proj=BA1F39653304A5B041B656915DC36B38&qid=0443DE
\BankTask{0443DE}{
Герметичный теплоизолированный сосуд разделили неподвижной перегородкой, способной проводить тепло, на две равные части. В первую часть сосуда поместили некоторое количества аргона при температуре \(328\,\text{К}\), а во вторую --- такое же количество аргона при температуре \(15\,^{\circ}\mathrm{C}\).
Теплоёмкость сосуда пренебрежимо мала.
Рассматриваются изменения, происходящие с аргоном при переходе к тепловому равновесию.
\begin{enumerate}
\item Температура газа во второй части сосуда повысилась.
\item Внутренняя энергия газа в первой части сосуда увеличилась.
\item При теплообмене газ в первой части сосуда отдавал теплоту, а газ во второй части сосуда её получал.
\item Через достаточно большой промежуток времени температура газов в обеих частях сосуда стала одинаковой и равной \(25\,^{\circ}\mathrm{C}\).
\item В результате теплообмена газ в первой части сосуда совершил положительную работу.
\end{enumerate}
}{
\smallskip\noindent\textbf{Решение от GPT}\par
Начальная температура во второй части \(T_2=15+273=288\,\text{К}\). Объёмы постоянны, поэтому работа каждого газа \(A=0\). По закону сохранения энергии для равных количеств одного газа
\[
T_{\text{равн}}=\frac{328+288}{2}=308\,\text{К}=35\,^{\circ}\mathrm{C}.
\]
\begin{enumerate}
\item Верно: температура второго газа повысилась от \(288\) до \(308\,\text{К}\).
\item Неверно: \(U=\frac32\nu RT\), поэтому при охлаждении первого газа его внутренняя энергия уменьшилась.
\item Верно: теплота передавалась от более горячего первого газа ко второму.
\item Неверно: равновесная температура равна \(35\,^{\circ}\mathrm{C}\).
\item Неверно: при неподвижной перегородке \(A=0\).
\end{enumerate}
Ответ: \(13\).
}

% Source 8054C4: https://ege.fipi.ru/bank/questions.php?search=1&proj=BA1F39653304A5B041B656915DC36B38&qid=8054C4
\BankTask{8054C4}{
Аргон массой $160\,\text{г}$ при комнатной температуре помещают в герметичный закрытый сосуд объёмом $20\,\text{л}$ и начинают нагревать. Масса газа в сосуде остаётся неизменной.

Установите соответствие между физическими величинами, характеризующими аргон, и формулами, выражающими их зависимость от абсолютной температуры $T$ газа в условиях данной задачи.

\begin{center}
\TaskMatchingColumns[.53\linewidth]{Физические величины}{
А) давление газа $p(T)$\par
Б) внутренняя энергия газа $U(T)$
}{Формулы}{
1) $aT$, $a=1662\,\dfrac{\text{Па}}{\text{К}}$\par
2) $bT$, $b=49{,}86\,\dfrac{\text{Дж}}{\text{К}}$\par
3) $\dfrac{c}{T}$, $c=80\,\text{Дж}\cdot\text{К}$\par
4) $dT$, $d=4{,}986\,\dfrac{\text{Па}}{\text{К}}$
}
\end{center}

\begin{center}
\begin{tabular}{|c|c|}\hline
А & Б \\ \hline
\phantom{00} & \phantom{00} \\ \hline
\end{tabular}
\end{center}
}{
\smallskip\noindent\textbf{Решение от GPT}\par
Используем молярную массу аргона $M=40\,\text{г/моль}$. Тогда $\nu=m/M=4\,\text{моль}$, $V=0{,}020\,\text{м}^3$.
\[
p(T)=\frac{\nu R}{V}T=\left(1662\,\frac{\text{Па}}{\text{К}}\right)T,
\qquad
U(T)=\frac32\nu RT=\left(49{,}86\,\frac{\text{Дж}}{\text{К}}\right)T.
\]
Давлению соответствует формула 1, внутренней энергии одноатомного газа — формула 2.

Ответ: А — 1, Б — 2.
}

% Source C427F7: https://ege.fipi.ru/bank/questions.php?search=1&proj=BA1F39653304A5B041B656915DC36B38&qid=C427F7
\BankTask{C427F7}{
Объём сосуда с идеальным газом уменьшили вдвое, выпустив половину газа и поддерживая температуру газа в сосуде постоянной. Как изменились при этом давление газа в сосуде и его внутренняя энергия?

Для каждой величины определите соответствующий характер изменения:

\begin{enumerate}
\item увеличилась
\item уменьшилась
\item не изменилась
\end{enumerate}


\TaskChangeTable{Давление газа в сосуде}{Внутренняя энергия газа в сосуде}
}{
\smallskip\noindent\textbf{Решение от GPT}\par
При постоянной температуре \(p=\nu RT/V\). Поскольку \(\nu_2=\nu_1/2\) и \(V_2=V_1/2\), получаем \(p_2=p_1\): давление не изменилось (\(3\)). Внутренняя энергия идеального газа при неизменной температуре пропорциональна количеству вещества: \(U_2=U_1/2\), поэтому она уменьшилась (\(2\)).

Ответ: \(32\).
}

% Source FB5B0D: https://ege.fipi.ru/bank/questions.php?search=1&proj=BA1F39653304A5B041B656915DC36B38&qid=FB5B0D
\BankTask{FB5B0D}{
В сосуде неизменного объёма при комнатной температуре находилась смесь неона и аргона, по \(2\,\text{моль}\) каждого. Половину смеси выпустили, а затем добавили в сосуд \(2\,\text{моль}\) аргона. Как в результате этого изменились парциальное давление аргона и давление смеси газов, если температура смеси оставалась неизменной? Для каждой величины определите соответствующий характер изменения:

\begin{enumerate}
\item увеличилась
\item уменьшилась
\item не изменилась
\end{enumerate}


\TaskChangeTable{Парциальное давление аргона}{Давление смеси газов в сосуде}
}{
\smallskip\noindent\textbf{Решение от GPT}\par
После выпуска половины смеси осталось по \(1\,\text{моль}\) каждого газа. После добавления аргона \(\nu_{\mathrm{Ar},2}=3\,\text{моль}\), а общее количество вещества \(\nu_2=4\,\text{моль}\). При постоянных \(T\) и \(V\) давление пропорционально количеству вещества. Поэтому парциальное давление аргона увеличилось в \(3/2\) раза (\(1\)), а давление смеси не изменилось, поскольку первоначально также было \(4\,\text{моль}\) газа (\(3\)).

Ответ: \(13\).
}

% Source 3BDF8F: https://ege.fipi.ru/bank/questions.php?search=1&proj=BA1F39653304A5B041B656915DC36B38&qid=3BDF8F
\BankTask{3BDF8F}{
В сосуде постоянного объёма абсолютную температуру одноатомного идеального газа увеличили в \(3\) раза, выпустив при этом \(\frac{1}{3}\) газа из сосуда. Как изменились в результате этого плотность газа в сосуде и его внутренняя энергия?

Для каждой величины определите соответствующий характер изменения:
\begin{enumerate}
\item увеличилась
\item уменьшилась
\item не изменилась
\end{enumerate}

\TaskChangeTable{Плотность газа в сосуде}{Внутренняя энергия газа в сосуде}
}{
\smallskip\noindent\textbf{Решение от GPT}\par
При постоянном объёме \(\rho=m/V\). После выпуска газа \(m_2=\frac{2}{3}m_1\), поэтому \(\rho_2=\frac{2}{3}\rho_1\): плотность уменьшилась (\(2\)).

Для одноатомного идеального газа \(U=\frac{3}{2}\nu RT\), поэтому
\[
\frac{U_2}{U_1}=\frac{\nu_2}{\nu_1}\frac{T_2}{T_1}=\frac{2}{3}\cdot3=2.
\]
Внутренняя энергия увеличилась (\(1\)). Ответ: \(21\).
}

% Source 2980: https://phys-ege.sdamgia.ru/problem?id=2980
\BankTask{2980}{
В сосуде с небольшой трещиной находится воздух. Воздух может медленно просачиваться сквозь трещину. Во время опыта объем сосуда уменьшили в $8$ раз, давление воздуха в сосуде увеличилось в $2$ раза, а его абсолютная температура увеличилась в $1{,}5$ раза. Во сколько раз изменилась внутренняя энергия воздуха в сосуде? (Воздух считать идеальным газом.)
}{
\smallskip\noindent\textbf{Решение}\par
Внутренняя энергия газа пропорциональна его температуре и количеству вещества (газа) в сосуде: $U= \frac{i}{2} \nu R T,$ где для воздуха можно принять $i=5.$ Согласно уравнению Клапейрона  — Менделеева, $pV=\nu RT.$ Следовательно, $U= \frac{i}{2} pV.$ Согласно условию задачи, p возросло в $2$ раза, а V уменьшилось в $8$ раз. Следовательно, U уменьшилась в $4$ раза.

Ответ: внутренняя энергия воздуха в сосуде уменьшилась в $4$ раза.
}

% Source 6939: https://phys-ege.sdamgia.ru/problem?id=6939
\BankTask{6939}{
Идеальный одноатомный газ медленно переводят из состояния $1$ в состояние $2$. Известно, что в процессе $1$ → $2$ давление газа изменялось прямо пропорционально его объему, а внутренняя энергия газа в этом процессе увеличилась на $6\,\text{Дж}$. Какую работу совершил газ в этом процессе?
}{
\smallskip\noindent\textbf{Решение}\par
По условию $p=const \cdot V.$ На диаграмме pV такой процесс изображается отрезком прямой, проходящей через начало координат. Работа равна площади фигуры под этим отрезком, то есть площади трапеции:

\[A= \frac{( p_2 + p_1 ) ( V_2 - V_1 )}{2} = \frac{( p_2V_2 - p_2V_1 + p_1V_2 - p_1V_1 )}{2} = \frac{( p_2V_2 - p_1V_1 ) + ( p_1V_2 - p_2V_1 )}{2} .\]

Первая скобка из уравнения состояния идеального газа ($pV=\nu RT$) равна $p_2V_2 - p_1V_1=\nu R ( T_2 - T_1 ) ,$ а вторая скобка равна $0$, так как по условию $\frac{p_1}{p_2} = \frac{V_1}{V_2} .$ Учитывая, что изменение внутренней энергии идеального одноатомного газа равно 

\[\Delta U=U_2 - U_1= \frac{3}{2} p_2V_2 - \frac{3}{2} p_1V_1,\]

получаем: $A= \frac{1}{3} \Delta U= \frac{1}{3} \cdot 6= 2\,\text{Дж}.$

Ответ: $2\,\text{Дж}$.
}

% Source 9042: https://phys-ege.sdamgia.ru/problem?id=9042
\BankTask{9042}{
В вакууме закреплен горизонтальный цилиндр (см. рис.). В цилиндре находится гелий, запертый поршнем. 
\begin{center}\begin{tikzpicture}[task graph,>=stealth]\draw[fill=gray!38](0,0)rectangle(5.7,1.85);\draw[fill=white](.8,.47)rectangle(5.7,1.38);\draw[fill=orange!35](.8,.47)rectangle(3.33,1.38);\fill(3.33,.47)rectangle(3.56,1.38);\draw[fill=gray!38](3.56,.47)--(4.13,.47)--(3.56,.72)--cycle;\draw[fill=gray!38](3.56,1.38)--(4.13,1.38)--(3.56,1.12)--cycle;\draw[->](6.6,.92)--(5.72,.92);\end{tikzpicture}\end{center}
Поршень массой $90\,\text{г}$ удерживается упорами и может скользить влево вдоль стенок цилиндра без трения. В поршень попадает пуля массой $10\,\text{г}$, летящая горизонтально со скоростью $400\,\text{м/с}$, и застревает в нем. Температура гелия в момент остановки поршня в крайнем левом положении возрастает на $64$ К. Чему равно количество вещества гелия в цилиндре? Считать, что за время движения поршня газ не успевает обменяться теплом с цилиндром и поршнем.
}{
\smallskip\noindent\textbf{Решение}\par
$1$.  Запишем закон сохранения импульса:

\[m v _1= ( m + M ) v _2,\]

где m и M  — масса пули и поршня соответственно, $v _1$  — скорость пули, $v _2$  — скорость поршня с застрявшей пулей.
Поршень будет двигаться со скоростью:

\[v _2= \frac{m v _1}{m + M} .\]

$2$.  Поршень с пулей будет обладать кинетической энергией:

\[E_{\,\text{кин}}= \frac{( m + M ) v _2 ^2}{2} = \frac{m ^2 v _1 ^2}{2 ( m + M )} .\]

Изменение кинетической энергии поршня будет равно работе газа $A=\Delta E_{\,\text{к}}=E_{\,\text{к}2} - E_{\,\text{к}1}= - \frac{m ^2 v _1 ^2}{2 ( m + M )} .$ Вследствие остановки поршня $E_{\,\text{к}2}=0.$
$3$.  Поскольку газ не успевает обменяться теплом с цилиндром и поршнем, то сжатие газа будет являться адиабатичным процессом. По первому началу термодинамики $- A=\Delta U.$ Вследствие этого процесса вся механическая энергия движения поршня с пулей пойдет на нагрев газа, поэтому:

\[\Delta U= \frac{3}{2} \nu R \Delta T= \frac{m ^2 v _1 ^2}{2 ( m + M )} .\]

$4$.  Найдем отсюда количество вещества в цилиндре:

\[\nu= \frac{m ^2 v _1 ^2}{3 ( m + M ) R \Delta T} = \frac{0{,}01 ^2 \cdot 400 ^2}{3 \cdot ( 0{,}01 + 0{,}09 ) \cdot 8{,}31 \cdot 64} \approx 0{,}1\,\text{моль}.\]

Ответ: $\nu\approx 0{,}1\,\text{моль}.$
}

% Source 7935: https://phys-ege.sdamgia.ru/problem?id=7935
\BankTask{7935}{
Два одинаковых теплоизолированных сосуда соединены короткой трубкой с краном. Объем каждого сосуда V  =  $1\,\text{м}^{3}$. В первом сосуде находится $\nu_{1}$  =  $1$ моль гелия при температуре $T_{1}$  =  $400$ К; во втором  — $\nu_{2}$  =  $3$ моль аргона при температуре $T_{2}$. Кран открывают. После установления равновесного состояния давление в сосудах р  =  $5{,}4$ кПа. Определите первоначальную температуру аргона $T_{2}$.
}{
\smallskip\noindent\textbf{Решение}\par
$1$.  Поскольку в указанном процессе газ не совершает работы и система является теплоизолированной, то в соответствии с первым законом термодинамики суммарная внутренняя энергия газов сохраняется:

\[\frac{3}{2} \nu_1RT_1 + \frac{3}{2} \nu_2RT_2= \frac{3}{2} ( \nu_1 + \nu_2 ) RT,\]

где T  — температура в объединенном сосуде в равновесном состоянии после открытия крана.
$2$.  В соответствии с уравнением Клапейрона  — Менделеева для конечного состояния можно записать: 

\[p ( 2V ) = ( \nu_1 + \nu_2 ) RT.\]

Исключая из двух записанных уравнений конечную температуру T, получаем искомое выражение для начальной температуры аргона:

\[T_2= \frac{2Vp}{\nu_2R} - \frac{\nu_1}{\nu_2} T_1= \frac{2 \cdot 1 \cdot 5{,}4 \cdot 10 ^3}{3 \cdot 8{,}31} - \frac{1}{3} \cdot 400 \approx 300\,\text{К}.\]

Ответ: $T_2 \approx 300\,\text{К}.$
}

\end{document}
